% !TEX root = master.tex
% !TEX program = lualatex
\chapter{Basics - 1}%
\label{cha:Basics - 1}

\section{C01a — Introduction, langage C et compilation}

\subsection{Objectif du bootcamp C}

Le langage C est introduit dans le cours de \textit{Computer Systems} comme un langage situé entre :

\begin{itemize}
    \item les langages haut niveau (Java, Scala)
    \item l’assembleur, très proche du processeur
\end{itemize}

Le C peut être vu comme un \textbf{assembleur structuré}, permettant :

\begin{itemize}
    \item une programmation proche de la machine
    \item un contrôle fin de la mémoire
    \item des programmes très efficaces
\end{itemize}

L’objectif du bootcamp est d’acquérir les bases du C, qui seront ensuite approfondies par les exercices et le projet du semestre.

\subsection{Public visé et avertissement important}

Le cours s’adresse à des étudiants sachant déjà programmer (typiquement en Java).

La syntaxe du C ressemble au Java car Java a été inspiré du C.
Cependant, la philosophie du langage est très différente.

\textbf{Attention :}
\begin{itemize}
    \item connaître Java ne signifie pas savoir programmer en C
    \item les tableaux se comportent différemment
    \item les chaînes ne sont pas des objets
    \item la gestion mémoire est manuelle
\end{itemize}

La pratique est donc indispensable pour maîtriser le langage.

\subsection{Nature du langage C}

Le C est :

\begin{itemize}
    \item un langage impératif
    \item non orienté objet
    \item compilé directement en code machine
    \item faiblement vérifié par le compilateur
    \item très ancien mais toujours très utilisé
\end{itemize}

Cette compilation directe en code machine permet d’obtenir des programmes très performants.

\subsection{Différences importantes avec Java}

\begin{itemize}
    \item pas d’orienté objet
    \item compilation en code machine (pas en bytecode)
    \item allocation mémoire explicite
    \item pas de garbage collector automatique
    \item accès direct à la mémoire via les pointeurs
    \item moins de vérifications du compilateur
\end{itemize}

Conséquences :

\begin{itemize}
    \item plus de contrôle
    \item plus de flexibilité
    \item mais plus de responsabilités
    \item et plus de bugs potentiels
\end{itemize}

\subsection{Structure générale d’un programme C}

Organisation typique :

\begin{enumerate}
    \item inclusions de bibliothèques
    \item variables globales (à éviter)
    \item fonctions auxiliaires
    \item fonction \texttt{main()}
\end{enumerate}

\subsection{La fonction \texttt{main}}

La fonction \texttt{main} est le point d’entrée du programme.

\begin{lstlisting}[language=C]
int main(void) {
    return 0;
}
\end{lstlisting}

\begin{itemize}
    \item retourne un entier au système d’exploitation
    \item 0 signifie succès
    \item une autre valeur représente un code d’erreur
\end{itemize}

\subsection{Exemple de programme}

Exemple : résolution d’une équation du second degré.

\begin{lstlisting}[language=C]
#include <stdio.h>
#include <math.h>

int main(void) {
    double b, c;
    printf("Enter b: ");
    scanf("%lf", &b);

    printf("Enter c: ");
    scanf("%lf", &c);

    double delta = b*b - 4*c;

    if (delta < 0) {
        puts("No real solution");
    } else if (delta == 0) {
        printf("One solution: %lf\n", -b/2);
    } else {
        printf("Two solutions\n");
    }

    return 0;
}
\end{lstlisting}

\subsection{Normes du langage C}

Versions importantes :

\begin{itemize}
    \item C89 : standard historique
    \item C99 : évolution majeure
    \item C11 / C17 : standards modernes
\end{itemize}

\subsubsection{Différences C89 vs C99}

\textbf{Commentaires}

C89 :
\begin{lstlisting}[language=C]
/* commentaire */
\end{lstlisting}

C99 :
\begin{lstlisting}[language=C]
// commentaire ligne unique
\end{lstlisting}

\textbf{Déclaration des variables}

C89 : toutes les variables doivent être déclarées au début du bloc.

C99 :
\begin{lstlisting}[language=C]
int a = 5;
printf("%d", a);
int b = 3;   // autorisé
\end{lstlisting}

Pour le cours, il est recommandé d’utiliser un standard récent :

\begin{lstlisting}[language=bash]
gcc -std=c17 program.c
\end{lstlisting}

\subsection{Compilation d’un programme}

Commande simple :

\begin{lstlisting}[language=bash]
gcc program.c -o program
\end{lstlisting}

Exécution :

\begin{lstlisting}[language=bash]
./program
\end{lstlisting}

\subsection{Étapes réelles de la compilation}

La compilation comporte plusieurs étapes :

\begin{enumerate}
    \item préprocesseur
    \item compilation en assembleur
    \item assemblage en fichier objet
    \item édition de liens (linking)
\end{enumerate}

\subsubsection{Options utiles}

\begin{itemize}
    \item \texttt{-E} : arrêt après préprocesseur
    \item \texttt{-S} : produire l’assembleur
    \item \texttt{-c} : produire fichier objet seulement
    \item \texttt{-o} : nom de l’exécutable
\end{itemize}

\subsection{Lien avec des bibliothèques}

Exemple avec la bibliothèque mathématique :

\begin{lstlisting}[language=bash]
gcc program.c -o program -lm
\end{lstlisting}

Sans \texttt{-lm}, le linker produit une erreur :

\begin{lstlisting}
undefined reference to sqrt
\end{lstlisting}

\subsection{Compile vs Build}

\begin{itemize}
    \item \textbf{Compile} : produit un fichier objet
    \item \textbf{Build} : compile + link = exécutable final
\end{itemize}

Pour les exercices : utiliser toujours le build complet.

\section{C01a — Basics 2 : Variables et types}

\subsection{Déclaration des variables}

En C, les variables ressemblent à Java syntaxiquement, mais :

\begin{itemize}
    \item ce ne sont pas des objets
    \item ce sont simplement des données en mémoire
\end{itemize}

Comme en Java, il est recommandé d’utiliser des noms explicites.

\subsubsection{Conventions de nommage}

Le C n’impose pas de convention stricte. Les plus courantes sont :

\begin{itemize}
    \item \textbf{camelCase} : \texttt{totalCount}
    \item \textbf{snake\_case} : \texttt{total\_count}
\end{itemize}

La convention dépend généralement du projet ou de l’équipe.

\subsection{Initialisation des variables}

\textbf{Point très important :}

En C, une variable n’est \textbf{pas initialisée par défaut}.

Une variable non initialisée possède une valeur indéterminée :
\[
\text{undefined behavior}
\]

\textbf{Règle : toujours initialiser les variables.}

\begin{lstlisting}[language=C]
int x;        // dangereux
int y = 0;    // correct
\end{lstlisting}

\subsection{Types disponibles}

Il y a moins de types qu’en Java et moins de bibliothèques standards.

\begin{itemize}
    \item pas de type \texttt{String}
    \item les chaînes seront traitées autrement
    \item pas de booléen en C89
\end{itemize}

Depuis C99, on peut utiliser :

\begin{lstlisting}[language=C]
#include <stdbool.h>
bool flag = true;
\end{lstlisting}

\subsection{Le mot-clé \texttt{const}}

\textbf{Important :} \texttt{const} ne signifie pas vraiment “constant”.

Il signifie :
\[
\text{read-only via this name}
\]

Autrement dit :

\begin{itemize}
    \item la variable ne peut pas être modifiée via ce nom
    \item mais sa valeur peut changer indirectement via la mémoire
\end{itemize}

\begin{lstlisting}[language=C]
const int a = 5;
a = 6;   // erreur compilation
\end{lstlisting}

Il est recommandé d’utiliser \texttt{const} chaque fois qu’une variable ne doit pas être modifiée.

\subsection{Sémantique de l’affectation}

\textbf{Différence fondamentale avec Java.}

\subsubsection{En Java (objets)}

\begin{lstlisting}[language=Java]
A a = new A();
A b = a;
b.x = 5;
\end{lstlisting}

Ici, \texttt{a} et \texttt{b} référencent le même objet.

Modifier \texttt{b} modifie aussi \texttt{a}.

C’est une \textbf{sémantique de référence}.

\subsubsection{En C}

\begin{lstlisting}[language=C]
int a = 5;
int b = a;
b = 10;
\end{lstlisting}

Ici :

\begin{itemize}
    \item \texttt{b} reçoit une copie de la valeur
    \item modifier \texttt{b} ne modifie pas \texttt{a}
\end{itemize}

C’est une \textbf{sémantique de valeur}.

\[
a = b \quad \Rightarrow \quad \text{copie mémoire}
\]

\subsection{Modificateurs de types}

On peut modifier les types de base :

\begin{itemize}
    \item \texttt{short}
    \item \texttt{long}
    \item \texttt{signed}
    \item \texttt{unsigned}
\end{itemize}

Exemples :

\begin{lstlisting}[language=C]
short int a;
long int b;
unsigned int c;
unsigned long int d;
\end{lstlisting}

Depuis C99 :

\begin{lstlisting}[language=C]
long long int x;
\end{lstlisting}

\subsection{Attention : tailles non fixées par la norme}

La norme C ne fixe pas la taille exacte des types.

Elle garantit seulement :

\[
\texttt{sizeof(char)} \leq \texttt{sizeof(short)} \leq \texttt{sizeof(int)} \leq \texttt{sizeof(long)}
\]

Donc :

\begin{itemize}
    \item ne jamais supposer qu’un \texttt{int} fait 32 bits
    \item attention à la portabilité
\end{itemize}

\subsection{Limites numériques}

Les limites des types sont définies dans :

\begin{lstlisting}[language=C]
#include <limits.h>
#include <float.h>
\end{lstlisting}

Exemples :

\begin{lstlisting}[language=C]
INT_MAX
LONG_MAX
ULONG_MAX
DBL_MAX
DBL_MIN
DBL_EPSILON
\end{lstlisting}

\subsection{Types à taille fixe (C99)}

Pour garantir un nombre exact de bits :

\begin{lstlisting}[language=C]
#include <stdint.h>
\end{lstlisting}

On peut utiliser :

\begin{lstlisting}[language=C]
int8_t
int16_t
int32_t
int64_t

uint8_t
uint16_t
uint32_t
uint64_t
\end{lstlisting}

Exemple :

\begin{lstlisting}[language=C]
uint8_t x = 255;   // exactement 8 bits
\end{lstlisting}

Ces types sont essentiels pour :

\begin{itemize}
    \item programmation système
    \item fichiers binaires
    \item réseau
    \item portabilité
\end{itemize}

\section{C01a — Basics 3 : Opérateurs et expressions}

\subsection{Remarque générale}

La plupart des opérateurs en C sont identiques à ceux de Java :

\begin{itemize}
    \item opérateurs arithmétiques
    \item opérateurs de comparaison
    \item opérateurs logiques
    \item priorités d’opérateurs
\end{itemize}

Cependant, quelques différences importantes doivent être retenues.

\subsection{Pas de concaténation de chaînes}

En C, il n’existe pas de type \texttt{String}.

Donc :

\begin{itemize}
    \item l’opérateur \texttt{+} ne concatène pas de chaînes
    \item les opérateurs arithmétiques sont strictement numériques
\end{itemize}

\subsection{Division entière vs division réelle}

Comme en Java, il existe deux types de division :

\begin{itemize}
    \item division entière : résultat entier
    \item division réelle : résultat flottant
\end{itemize}

\begin{lstlisting}[language=C]
int a = 3/2;      // vaut 1
double b = 3/2;   // vaut 1.0
double c = 3.0/2; // vaut 1.5
\end{lstlisting}

\textbf{Erreur classique :} calculer une moyenne avec des entiers.

\subsection{Toute expression a une valeur}

En C, toute expression :

\begin{itemize}
    \item produit un effet
    \item retourne une valeur
\end{itemize}

Par exemple :

\begin{lstlisting}[language=C]
x = 3;
\end{lstlisting}

Cette expression :

\begin{itemize}
    \item affecte 3 à \texttt{x}
    \item retourne la valeur 3
\end{itemize}

\subsection{Erreur classique : affectation dans un if}

Le code suivant compile :

\begin{lstlisting}[language=C]
if (x = 3) {
    ...
}
\end{lstlisting}

Pourquoi ?

\begin{itemize}
    \item \texttt{x = 3} vaut 3
    \item 3 est différent de 0
    \item donc la condition est toujours vraie
\end{itemize}

\textbf{C’est une erreur très fréquente en C.}

\subsubsection{Rappel}

\begin{itemize}
    \item \texttt{==} teste l’égalité
    \item \texttt{=} affecte une valeur
\end{itemize}

\subsubsection{Astuce pour éviter l’erreur}

Écrire les constantes à gauche :

\begin{lstlisting}[language=C]
if (3 == x)   // correct
\end{lstlisting}

Car :

\begin{lstlisting}[language=C]
if (3 = x)    // erreur compilation
\end{lstlisting}

\subsection{Valeurs de vérité en C}

En C :

\begin{itemize}
    \item toute valeur non nulle est vraie
    \item seule la valeur 0 est fausse
\end{itemize}

Donc :

\begin{lstlisting}[language=C]
if (5)      // vrai
if (-2.4)   // vrai
if (0)      // faux
\end{lstlisting}

\subsection{Expressions logiques}

Toute expression peut être utilisée comme condition.

\begin{lstlisting}[language=C]
int x = 10;
if (x) { ... }     // vrai car x ≠ 0
\end{lstlisting}

\subsection{Opérateur virgule (avancé)}

Le C possède un opérateur absent en Java : l’opérateur virgule.

Syntaxe :

\[
E_1 , E_2
\]

Effet :

\begin{itemize}
    \item évalue d’abord \(E_1\)
    \item puis évalue \(E_2\)
    \item la valeur de l’expression est celle de \(E_2\)
\end{itemize}

\subsubsection{Exemple simple}

\begin{lstlisting}[language=C]
x = (3,4);
\end{lstlisting}

\begin{itemize}
    \item 3 est évalué
    \item 4 est évalué
    \item l’expression vaut 4
\end{itemize}

Donc :

\[
x = 4
\]

\subsubsection{Exemple avec affectation}

\begin{lstlisting}[language=C]
x = (a = b, 3*a);
\end{lstlisting}

Étapes :

\begin{enumerate}
    \item on copie \texttt{b} dans \texttt{a}
    \item on calcule \texttt{3*a}
    \item cette valeur est affectée à \texttt{x}
\end{enumerate}

C’est équivalent à :

\begin{lstlisting}[language=C]
a = b;
x = 3*a;
\end{lstlisting}

\subsubsection{Priorité avec l’affectation}

Attention :

\begin{lstlisting}[language=C]
x = 3,4; // pas pour la déclaration !
\end{lstlisting}

est interprété comme :

\begin{lstlisting}[language=C]
(x = 3), 4;
\end{lstlisting}

Donc :

\begin{itemize}
    \item \texttt{x} reçoit 3
    \item l’expression complète vaut 4
\end{itemize}

\subsection{Conseil de style}

L’opérateur virgule rend le code difficile à lire.

\textbf{Recommandation :}
\begin{itemize}
    \item éviter de l’utiliser
    \item mais savoir le lire
\end{itemize}

\subsection{Opérateurs bit à bit}

Les opérateurs bit à bit sont les mêmes qu’en Java :

\begin{itemize}
    \item \texttt{\&} AND
    \item \texttt{|} OR
    \item \texttt{\^} XOR
    \item \texttt{\textasciitilde} NOT
    \item \texttt{<<} décalage gauche
    \item \texttt{>>} décalage droite
\end{itemize}

Ils sont importants pour la programmation système.

\section{C01a — Basics 4 : Structures de contrôle}

\subsection{Rôle des structures de contrôle}

Les structures de contrôle permettent de modifier l’exécution linéaire d’un programme :

\begin{itemize}
    \item branchements conditionnels : \texttt{if}
    \item boucles conditionnelles : \texttt{while}, \texttt{do while}
    \item itérations : \texttt{for}
\end{itemize}

Elles contrôlent toujours un \textbf{bloc d’instructions}, défini par des accolades :

\begin{lstlisting}[language=C]
if (cond) {
    ...
}
\end{lstlisting}

Un bloc peut contenir ses propres variables locales.

\subsection{Portée des variables}

La portée d’une variable correspond à la zone du programme où elle est accessible.

On distingue :

\begin{itemize}
    \item \textbf{portée locale} : variable définie dans un bloc
    \item \textbf{portée globale} : variable définie hors de tout bloc
\end{itemize}

Exemple :

\begin{lstlisting}[language=C]
if (cond) {
    int y = 3;
}
// y n’existe plus ici
\end{lstlisting}

\subsection{Masquage de variables}

Contrairement à Java, le C autorise de redéclarer une variable dans un bloc interne.

\begin{lstlisting}[language=C]
int y = 10;

{
    int y = 1;   // masque la variable extérieure
    printf("%d", y);   // affiche 1
}

printf("%d", y);       // affiche 10
\end{lstlisting}

La règle est :

\[
\textbf{on utilise toujours la variable de portée la plus proche}
\]

Conseil : éviter de réutiliser les mêmes noms.

\subsection{Variables globales}

Une variable globale est définie hors de toute fonction.

\begin{lstlisting}[language=C]
int max = 5;
\end{lstlisting}

Problèmes :

\begin{itemize}
    \item accessible depuis tout le programme
    \item difficile à contrôler
    \item effets de bord possibles
\end{itemize}

\textbf{Conseil : éviter les variables globales}
(sauf constantes globales).

\subsection{Le if}

Identique à Java :

\begin{lstlisting}[language=C]
if (cond) {
    ...
} else {
    ...
}
\end{lstlisting}

\subsubsection{Conseil important}

Toujours mettre des accolades, même pour une seule instruction.

\begin{lstlisting}[language=C]
if (cond)
    do_something();   // dangereux
\end{lstlisting}

Si on ajoute une ligne ensuite, elle ne sera pas dans le bloc.

Cette erreur a déjà causé des bugs industriels majeurs.

\subsection{else if}

Il n’existe pas réellement en C :  
c’est simplement un \texttt{if} dans un \texttt{else}.

\begin{lstlisting}[language=C]
if (a==1) { ... }
else if (a==2) { ... }
else { ... }
\end{lstlisting}

\subsection{switch}

Utilisé pour tester plusieurs valeurs possibles :

\begin{lstlisting}[language=C]
switch(x) {
    case 1:
        ...
        break;
    case 2:
        ...
        break;
    default:
        ...
}
\end{lstlisting}

\subsection{Boucle while}

Boucle conditionnelle a priori :

\begin{lstlisting}[language=C]
while (cond) {
    ...
}
\end{lstlisting}

\subsection{Boucle do while}

Boucle conditionnelle a posteriori  
(le bloc s’exécute au moins une fois).

\begin{lstlisting}[language=C]
do {
    ...
} while (cond);

\end{lstlisting}

\subsection{Boucle for}

Forme générale :

\begin{lstlisting}[language=C]
for (initialisation ; condition ; mise_a_jour) {
    ...
}
\end{lstlisting}

Équivalent à :

\begin{lstlisting}[language=C]
initialisation;
while (condition) {
    ...
    mise_a_jour;
}
\end{lstlisting}

\subsubsection{Portée}

Les variables déclarées dans le \texttt{for} sont locales à la boucle :

\begin{lstlisting}[language=C]
for (int i=0; i<5; i++) { ... }
// i n’existe plus ici
\end{lstlisting}

\subsubsection{Conseil}

Utiliser \texttt{for} pour les itérations connues :

\begin{itemize}
    \item plus lisible
    \item plus compact
    \item plus clair
\end{itemize}

\subsection{Initialisation multiple dans for}

On peut déclarer plusieurs variables :

\begin{lstlisting}[language=C]
for (int i=0, s=0; i<5; i++) {
    s += i;
}
\end{lstlisting}

\subsection{Mise à jour complexe}

On peut utiliser l’opérateur virgule :

\begin{lstlisting}[language=C]
for (int i=0, s=0; i<5; i++, s+=i) {
    ...
}
\end{lstlisting}

Mais ce style est peu lisible.

\textbf{Conseil : éviter ce genre d’écriture.}

\subsection{break et continue}

Toutes les structures de contrôle peuvent être interrompues par :

\begin{itemize}
    \item \texttt{break} : sortie de la boucle
    \item \texttt{continue} : passe à l’itération suivante
\end{itemize}

\subsubsection{Conseil}

Éviter leur usage excessif :

\begin{itemize}
    \item rend le code difficile à suivre
    \item peut introduire des bugs subtils
    \item donne un style dit \textit{spaghetti code}
\end{itemize}

Préférer écrire une condition claire plutôt que d’utiliser
\texttt{break} ou \texttt{continue}.

\section{C01a — Basics 5 : Fonctions}

\subsection{Pourquoi utiliser des fonctions}

Les fonctions servent principalement à deux choses :

\begin{itemize}
    \item mettre en évidence un traitement important
    \item réutiliser du code sans duplication
\end{itemize}

\textbf{Règle fondamentale : ne jamais copier-coller du code.}

Si une modification conceptuelle nécessite plusieurs modifications dans le code,
cela signifie qu’il y a duplication.

La solution consiste à factoriser ce traitement dans une fonction.

\subsection{Fonctions en C}

Les fonctions en C ressemblent à celles de Java, mais :

\begin{itemize}
    \item il n’y a pas d’objets
    \item il n’y a pas de méthodes
    \item toutes les fonctions sont globales
\end{itemize}

\subsection{Prototype d’une fonction}

Le C introduit une notion absente de Java : le \textbf{prototype}.

Le prototype correspond à l’en-tête de la fonction terminé par un point-virgule.

Exemple :

\begin{lstlisting}[language=C]
double moyenne(double a, double b);
\end{lstlisting}

Le prototype sert à annoncer l’existence de la fonction avant son utilisation.

\textbf{Règle : une fonction doit être déclarée avant d’être appelée.}

\subsection{Prototype vs définition}

La définition contient déjà le prototype :

\begin{lstlisting}[language=C]
double moyenne(double a, double b) {
    return (a+b)/2;
}
\end{lstlisting}

Donc :

\begin{itemize}
    \item soit on met le prototype avant
    \item soit on met directement la définition avant
\end{itemize}

Dans les programmes simples à un fichier, on peut simplement placer les
définitions avant \texttt{main}.

Les prototypes deviennent essentiels dans les projets multi-fichiers.

\subsection{Arguments dans les prototypes}

Dans un prototype, seuls sont nécessaires :

\begin{itemize}
    \item le type de retour
    \item le nom de la fonction
    \item les types des arguments
\end{itemize}

Les noms des arguments ne sont pas obligatoires :

\begin{lstlisting}[language=C]
double moyenne(double, double);
\end{lstlisting}

Cependant, il est recommandé d’indiquer des noms pour la lisibilité.

\subsection{Fonctions sans arguments}

En C, une fonction sans argument doit explicitement indiquer \texttt{void} :

\begin{lstlisting}[language=C]
int f(void);
\end{lstlisting}

Contrairement à Java :

\begin{lstlisting}[language=Java]
int f()    // OK en Java
\end{lstlisting}

En C, écrire :

\begin{lstlisting}[language=C]
int f()
\end{lstlisting}

correspond à un ancien style et ne doit pas être utilisé.

\subsection{Fonctions sans valeur de retour}

On utilise le type \texttt{void} :

\begin{lstlisting}[language=C]
void afficher(int x) {
    printf("%d", x);
}
\end{lstlisting}

\subsection{Fonctions qui ne retournent jamais}

Depuis C11, on peut utiliser :

\begin{lstlisting}[language=C]
_Noreturn void f(void);
\end{lstlisting}

Ces fonctions quittent le programme définitivement (usage système).

\subsection{Passage des arguments}

En programmation, il existe deux modèles :

\begin{itemize}
    \item passage par valeur
    \item passage par référence
\end{itemize}

\subsection{Passage par valeur en C}

Le C ne possède que le passage par valeur.

Exemple :

\begin{lstlisting}[language=C]
void f(int x) {
    x++;
}

int main(void) {
    int val = 1;
    f(val);
    printf("%d", val);
}
\end{lstlisting}

Résultat :

\[
val = 1
\]

Car \texttt{x} est une copie.

\subsection{Simuler un passage par référence}

Pour modifier une variable appelante, on utilise un pointeur.

Conceptuellement :

\begin{itemize}
    \item on passe la valeur de l’adresse
    \item donc un pointeur
\end{itemize}

Deux choses à faire :

\begin{enumerate}
    \item utiliser \texttt{*} dans la fonction \textrightarrow crée un pointeur
    \item utiliser \texttt{\&} à l’appel = adresse de 
\end{enumerate}

\subsubsection{Exemple}

\begin{lstlisting}[language=C]
void f(int *x) {
    (*x)++;
}

int main(void) {
    int val = 1;
    f(&val);
    printf("%d", val);
}
\end{lstlisting}

Résultat :

\[
val = 2
\]

On parle couramment de passage par référence,
mais techncquement c’est un passage par valeur de pointeur.

\subsection{Surcharge des fonctions}

Contrairement à Java, le C ne supporte pas la surcharge.

Impossible d’écrire :

\begin{lstlisting}[language=C]
void affiche(int x);
void affiche(double x);
\end{lstlisting}

Il faut utiliser des noms différents :

\begin{lstlisting}[language=C]
void affiche_int(int x);
void affiche_double(double x);
\end{lstlisting}

